\chapter{Einleitung}
\label{chap:einleitung}

Die fortschreitende Digitalisierung, Vernetzung und der massive Ausbau moderner Internet- und Cloud-Infrastrukturen gehen Hand in Hand mit einer stetig wachsenden Bedrohungslandschaft im Bereich der Cybersicherheit. Unter den zahlreichen Angriffsszenarien stellen Distributed-Denial-of-Service-Angriffe (DDoS) nach wie vor eine der destruktivsten und am schwersten abzuwehrenden Gefahren dar. Durch den koordinierten Einsatz verteilter Botnets versuchen Angreifer, kritische Ressourcen, Server und Netzwerkinfrastrukturen durch massiven Datenverkehr zu überlasten und somit für legitime Nutzer unzugänglich zu machen \cite{mirkovicTaxonomyDDoSAttack2004}.

Traditionelle Erkennungsmechanismen, die auf statischen Signaturen oder einfachen Schwellenwerten basieren, stoßen in hochdynamischen modernen Netzwerken zunehmend an ihre Leistungsgrenzen. Moderne Angriffe zeichnen sich durch polymorphes Verhalten, verschleierte Traffic-Muster und die Nachahmung legitimen Nutzerverhaltens aus. Vor diesem Hintergrund gewinnt der Einsatz von Verfahren der Künstlichen Intelligenz (KI) und des maschinellen Lernens in der Netzwerkabwehr drastisch an Bedeutung \cite{kumarComprehensiveReviewVulnerabilities2024}. Durch die automatisierte Analyse komplexer Datenströme und die Erkennung nicht-linearer Anomalien versprechen KI-gestützte Ansätze eine signifikante Steigerung der Robustheit und Reaktionsgeschwindigkeit.

\section{Zielsetzung und Forschungsfragen}
Das Hauptziel dieser Bachelorarbeit besteht darin, die Leistungsfähigkeit und Robustheit von KI-basierten Klassifikationsverfahren zur Erkennung von DDoS-Angriffen systematisch zu untersuchen und zu bewerten. Im Fokus stehen dabei insbesondere die Validierung anhand aktueller Evaluationsdatensätze sowie die Analyse des Einflusses von Störungen auf die Detektionsgenauigkeit.

Dabei werden folgende zentrale Forschungsfragen adressiert:
\begin{itemize}
    \item Welche aktuellen Klassifikationsverfahren und Deep-Learning-Architekturen eignen sich am besten für die Identifikation von DDoS-Angriffen in modernen Netzwerktopologien?
    \item Wie wirken sich stochastische Kanalstörungen und Datenvariationen auf die Robustheit der Modelle aus?
    \item Welche Performance-Unterschiede zeigen sich bei der Evaluierung auf standardisierten Benchmark-Datensätzen?
\end{itemize}

\section{Aufbau der Arbeit}
Die vorliegende Arbeit gliedert sich in folgende Kapitel:
\begin{itemize}
    \item \textbf{Kapitel 2 (Grundlagen):} Vermittelt die theoretischen Grundlagen zu DDoS-Angriffsvektoren, Taxonomien und den verwendeten Klassifikationsverfahren.
    \item \textbf{Kapitel 3 (Verwandte Arbeiten / Related Work):} Gibt einen Überblick über den aktuellen Forschungsstand und bestehende Lösungsansätze in der Literatur.
    \item \textbf{Kapitel 4 (Datensätze \& Evaluierung):} Beleuchtet die verwendeten Evaluationsdatensätze und deren Eigenschaften.
    \item \textbf{Kapitel 5 (Implementierung):} Beschreibt die technische Umsetzung und Modellarchitektur.
    \item \textbf{Kapitel 6 (Evaluation):} Präsentiert und diskutiert die experimentellen Ergebnisse.
    \item \textbf{Kapitel 7 (Fazit):} Fasst die Ergebnisse zusammen und gibt einen Ausblick auf zukünftige Forschungsarbeiten.
\end{itemize}